\section{Background and Related Work}
\label{sec:related-work}

\subsection{Concept Erasure and Its Reference}

Post-training concept erasure aims to suppress a model's knowledge of a target
concept while preserving behavior unrelated to that concept. More broadly,
machine unlearning is commonly framed as approximating the model that would
have resulted had the data to be forgotten not contributed to training
\citep{wichert2024rethinking}. This framing suggests two distinct evaluation
questions. First, does an intervention reduce target knowledge while
preserving other capabilities? Second, does the resulting model resemble one
trained without the target learning signal?

Common evaluations compare an edited model with its original version, measuring
target suppression and capability preservation but not similarity to a model
trained without the target data \citep{yao2024machineunlearning,li2024wmdp}.
A matched training intervention is therefore needed to evaluate the latter.

\subsection{Concept Exclusion with LMEnt}

LMEnt provides a Wikipedia-based pretraining corpus with entity annotations,
an index for retrieving chunks associated with particular entities, and
language models trained on the corpus \citep{gottesman2025lment}. Its entity
index makes it possible to define a concept-level training intervention:
chunks linked to a selected set of entities can be prevented from contributing
to the training loss.

We refer to a model trained under such an intervention as a
\emph{concept-excluded twin}. A twin is trained under otherwise matched
conditions but receives no loss on the selected concept-associated chunks. It
therefore provides a reference for the behavioral effect of excluding a
defined portion of the training signal. This reference is more informative
than the original model alone when asking whether post-training erasure
reproduces the outcome of concept exclusion.

Concept exclusion does not guarantee that a model contains no target
knowledge. The entity-based exclusion set may miss indirect or unlinked
references while also covering text about neighboring topics. Consequently,
a twin may retain target knowledge and differ from the full model outside the
target concept. It therefore represents the effect of the specified
loss-exclusion procedure, not a perfectly concept-free model.

\subsection{Post-Training Concept-Erasure Methods}

We study three post-training methods that intervene at different locations in
the model. EMBER edits input-token embeddings rather than internal transformer
layers \citep{suslik2026ember}. It applies sparse matrix factorization to
embeddings obtained from target and neutral sentences, identifies factors
associated with the target concept, and subtracts scaled contributions of
those factors from selected token embeddings. The remaining model parameters
are left unchanged.

Representation Misdirection for Unlearning (RMU) modifies internal model
representations \citep{li2024wmdp}. It fine-tunes selected transformer weights
so that representations of forget-set inputs approach a random control
direction, while a retain objective encourages representations on non-target
inputs to remain close to those of the original model. RMU therefore combines
a target-directed intervention with an explicit preservation objective.

Semi-nonnegative matrix factorization (SNMF) identifies interpretable features
from groups of co-activated MLP neurons \citep{shafran2026snmf}. In the
erasure adaptation evaluated by EMBER, concept-associated features are mapped
to directions that can be removed through edits to MLP weights
\citep{suslik2026ember}. Because EMBER acts on token embeddings whereas RMU
and the SNMF-based method act on internal weights, EMBER can also be combined
with either MLP-level method in a staged intervention.

\subsection{Evaluation of Concept Erasure}

The EMBER study evaluated target-concept efficacy, specificity on related
domains, preservation on general benchmarks, output coherence, and robustness
to subsequent relearning \citep{suslik2026ember}. It assessed both
multiple-choice and open-ended target knowledge and compared EMBER with RMU,
SNMF-based MLP erasure, and staged combinations of these methods. Its
normalized evaluation used the original pre-erasure model as the primary
reference.

RMU was originally evaluated by measuring suppression on the WMDP benchmark
alongside capability retention on MMLU and MT-Bench
\citep{li2024wmdp}. These evaluations establish whether an intervention
suppresses selected behavior while retaining broader capabilities. They do
not, however, establish whether the edited model resembles a matched model
trained without the relevant learning signal.

Our evaluation separates these objectives. We measure erasure efficacy and
preservation relative to the full model, while independently measuring
similarity to a concept-excluded twin. The latter comparison uses
correct-answer negative log-likelihood and teacher-forced full-vocabulary KL
divergence on the same sentence completions. This design tests whether target
suppression coincides with the behavioral effects of concept exclusion during
training.
